\section{Training Recipe and Evidence Inventory}
\label{app:recipe}

Table~\ref{tab:recipe} summarizes the resolved run arguments. The base width
is 384 with six heads and MLP expansion four. Stochastic depth is configured
at 0.1; \mn{} constructs separate local and latent ramps, unlike a single
twelve-block dense ramp. The global and flattened routing controls use their
dense/historical scoring paths, while radius-2/3 routing uses the sparse
spatial backend; the routing support and its numerical implementation should
therefore both be recorded. EMA validation replaces raw-model metrics before
the summary row is written, so \texttt{eval\_top1} is an EMA metric.

\begin{table}[h]
\centering\small
\begin{tabular}{lll}
\toprule
Setting & 224-pixel scratch & 384/512-pixel fine-tune\\
\midrule
Epochs & 150 or 300 & 30 at 384; 15 at 512\\
Effective global batch & 1024 & 1024\\
Optimizer & AdamW & AdamW\\
Learning rate & $5\!\times\!10^{-4}$ or $7.5\!\times\!10^{-4}$ & $10^{-5}$\\
Minimum learning rate & $10^{-5}$ & $10^{-6}$\\
Weight decay & 0.05 & $10^{-8}$\\
Warmup epochs & 5 (probe: 20) & 1\\
Warmup learning rate & $10^{-6}$ & $10^{-6}$\\
EMA decay & 0.99996 & 0.9996\\
Mixup / CutMix & 0.8 / 1.0 & 0 / 0\\
Label smoothing & 0.1 & 0.1\\
Random erasing & 0.25 & 0.25\\
RandAugment & \texttt{rand-m9-mstd0.5-inc1} & same\\
Gradient clip norm & 1.0 & 1.0\\
Validation crop percentage & 0.9 & 1.0\\
\mn{} auxiliary-logit weight & ramp to 0.05 & 0.05\\
Teacher / feature / routing distillation & disabled & disabled\\
\bottomrule
\end{tabular}
\caption{Main recipe. Launch records take precedence over configuration-file
defaults, especially for microbatch size and gradient accumulation.}
\label{tab:recipe}
\end{table}

\paragraph{Classifier implementation.}
The main configuration uses patch size 8, width 384, six heads, six local
blocks, and six latent blocks. At 224 pixels this gives 784 patch slots; the
local attention band extends 16 sequence positions on either side of a patch
query, while the CLS query attends globally. Spatial routing uses radius 3,
a $384\to64$ metric projection, similarity temperature 1, and input-gradient
scale 0.1 for the metric. At compression factor 2, the six soft donor budgets
are $(65,65,65,65,65,67)$ and the hard gather retains 392 patches. There is
no further merging in the latent encoder. Training partitions donors and
receivers randomly for each image; evaluation alternates deterministic
odd/even partitions across routing steps.

\paragraph{Selection and tracing.}
The threshold-based soft top-$k$ uses scale 30 and normalizes row energies by
their maximum-minus-minimum range across the current minibatch. Thus soft
selection can depend on batch companions; the released operator has no guard
for a zero range. The final hard-gather indices are detached and selected
carriers are sorted by a detached, mass-weighted center of their flattened
source indices. The default center-tracing mode does not retain a full source
membership matrix. Spatial scoring uses sparse neighbor lists, although
model initialization constructs a dense Boolean adjacency buffer.

During training, an auxiliary readout averages full-grid features using
soft-selection weights, applies the shared classification head, and adds its
weighted logits to the latent CLS logits. It is absent at evaluation and is
not a separate auxiliary loss. The 150-epoch \mn{} runs start this
contribution at epoch 10
and ramp it over 10 epochs; the 300-epoch run uses epoch 20 and a 20-epoch
ramp. Fine-tunes enable weight 0.05 from the start. The
recorded environment is Python 3.10.20, PyTorch 2.7.1+cu118, torchvision
0.22.1+cu118, timm 0.9.11, FlashAttention 2.7.3 and Triton 3.3.1 on H20-3e
GPUs. Both four and eight ranks divide 50,000, avoiding validation-sampler
padding for this split. Other world sizes require a separate check.

Three nominal launcher failures occurred after complete metric histories:
the 150-epoch DeiT-S/16, the compound $\lambda=4$ run, and dense 512-pixel
fine-tuning. Complete summaries and supplied checkpoint receipts distinguish
them from early failures. Receipts record checks on another machine, not
weights independently loaded in this audit. The $\lambda$ curriculum and
512-pixel \mn{} fine-tune are now complete 150-epoch and 15-epoch records.
The weight-decay probe was abandoned at 140/150. The DeiT-initialized and
progressive-latent runs are complete 150-epoch records.

\begin{table}[h]
\centering\footnotesize
\begin{tabular}{lrlr}
\toprule
Run ID & Epochs & Status & Best top-1 \\
\midrule
\texttt{s2\_deit\_p16\_150e} & 150/150 & Complete & 77.196 \\
\texttt{s2\_deit\_p8\_150e} & 150/150 & Complete & 79.992 \\
\texttt{s2\_mn\_flat\_w8\_150e} & 150/150 & Complete & 78.290 \\
\texttt{s2\_mn\_global\_150e} & 150/150 & Complete & 77.218 \\
\texttt{s2\_mn\_l4\_150e} & 150/150 & Complete & 78.366 \\
\texttt{s2\_mn\_r2\_150e} & 150/150 & Complete & 78.942 \\
\texttt{s2\_mn\_r3\_150e} & 150/150 & Complete & 78.888 \\
\texttt{s2\_mn\_r3\_lr75\_150e} & 150/150 & Complete & 79.316 \\
\texttt{s3\_deit\_p16\_300e} & 300/300 & Complete & 79.692 \\
\texttt{s3\_deit\_p8\_300e} & 300/300 & Complete & 82.246 \\
\texttt{s3\_mn\_r3\_lr75\_300e} & 300/300 & Complete & 81.368 \\
\texttt{s3\_deit\_p8\_ft384\_30e} & 30/30 & Complete & 82.976 \\
\texttt{s3\_mn\_r3\_ft384\_30e} & 30/30 & Complete & 82.582 \\
\texttt{s3\_mn\_r5\_ft384\_30e} & 30/30 & Complete & 82.574 \\
\texttt{s4\_mn\_lamcurr\_150e} & 150/150 & Complete & 79.310 \\
\texttt{s4\_mn\_warm20\_150e} & 150/150 & Complete & 79.238 \\
\texttt{s4\_deit\_ft512\_15e} & 15/15 & Complete & 83.068 \\
\texttt{s4\_mn\_ft512\_15e} & 15/15 & Complete & 82.618 \\
\texttt{s5\_mn\_dp07\_150e} & 150/150 & Complete & 79.534 \\
\texttt{s5\_mn\_lr1e3\_150e} & 150/150 & Complete & 79.350 \\
\texttt{s5\_mn\_wd03\_150e} & 140/150 & Partial & 79.054 \\
\texttt{s6\_deitB\_p16\_300e} & 300/300 & Complete & 80.442 \\
\texttt{s7\_mn\_deitinit\_150e} & 150/150 & Complete & 80.574 \\
\texttt{s7\_mn\_proglatent\_150e} & 150/150 & Complete & 79.392 \\
\bottomrule
\end{tabular}

\caption{Status through the September 23 packet. The abandoned weight-decay
metric is included for accounting only and is not an endpoint. The DeiT-B best score is a
pre-overfitting peak (Appendix~\ref{app:deitb}). The companion data retain
full precision and source summary lines.}
\label{tab:runs}
\end{table}
\FloatBarrier

\subsection{Initialization and Latent-Compression Diagnostics}

Table~\ref{tab:stage7-diagnostics} uses the completed drop-path-0.07 run as
the 150-epoch scratch control. Initializing \mn{} from the 300-epoch
DeiT-S/8 checkpoint raises best EMA top-1 by 1.040 points. This result shows
initialization sensitivity, not equal-total-compute superiority: the source
representation already consumed 300 training epochs, the experiment has one
seed, and its 80.574\% endpoint remains below both the 300-epoch scratch
\mn{} result (81.368\%) and its dense source (82.246\%). The original
first-launch loading log was not delivered after a host migration; the
record retains a contemporaneous 152-loaded/25-missing receipt and the early
learning curve, so we treat this as a diagnostic rather than a primary causal
result.

The progressive-latent variant adds 192 latent-stage merges and reaches
79.392\%, 0.142 points below the same control. It changes the final configured
budget from approximately 392 to 200 patches, so it is an
accuracy--compression probe rather than a matched-budget test of gradual
versus abrupt scheduling. With one seed, the result establishes no observed
gain but does not establish a statistically resolved degradation. Training-log
throughput is not used as deployment latency.

\begin{table}[h]
\caption{Completed 150-epoch Stage-7 diagnostics, all with seed 42 and
drop-path 0.07. Deltas use the scratch control in the first row.}
\label{tab:stage7-diagnostics}
\centering\small
\begin{tabular}{llrr}
\toprule
Variant & Setting / change from control & Best (\%) & $\Delta$ (pp) \\
\midrule
Scratch control & drop-path 0.07 & 79.534 & 0.000 \\
DeiT initialization & initialize from DeiT-S/8 300e & 80.574 & +1.040 \\
Progressive latent & 192 additional latent merges & 79.392 & -0.142 \\
\bottomrule
\end{tabular}

\end{table}

\section{Detailed Budget and Timing Comparisons}
\label{app:latency}

Table~\ref{tab:sweepdetail} gives the values behind the main accuracy panels.
The \mn{} checkpoint uses radius 3 at 224 pixels and radius 5.143 after
384-pixel fine-tuning. Final patch budgets match within a row, but layerwise
attention and MLP inputs do not. ToMe and PiToMe shorten sequences within
successive blocks, while \mn{} gathers once between stages. The companion
CSV also contains the two resolution-transfer panels using 224-pixel
checkpoints without fine-tuning.

\begin{table}[h]
\centering\small
\begin{tabular}{lrrrrr}
\toprule
Checkpoint / pixels & Patches & MergeNet & ToMe & PiToMe & $\Delta$ ToMe \\
\midrule
224 trained / 224 & 523 & 81.380 & 82.038 & 81.876 & +0.658 \\
224 trained / 224 & 392 & 81.360 & 81.932 & 81.742 & +0.572 \\
224 trained / 224 & 314 & 81.074 & 81.814 & 81.560 & +0.740 \\
224 trained / 224 & 262 & 80.850 & 81.830 & 81.500 & +0.980 \\
224 trained / 224 & 196 & 80.228 & 81.752 & 81.206 & +1.524 \\
384 tuned / 384 & 1537 & 82.500 & 82.868 & 82.658 & +0.368 \\
384 tuned / 384 & 1152 & 82.562 & 82.866 & 82.534 & +0.304 \\
384 tuned / 384 & 923 & 82.230 & 82.800 & 82.452 & +0.570 \\
384 tuned / 384 & 770 & 82.002 & 82.792 & 82.334 & +0.790 \\
384 tuned / 384 & 576 & 81.300 & 82.742 & 82.296 & +1.442 \\
\bottomrule
\end{tabular}

\caption{Top-1 percentages at equal final patch counts; $\Delta$ is ToMe minus
\mn{} in percentage points. These values cannot be paired with
Table~\ref{tab:latency384} to construct an accuracy--latency frontier.}
\label{tab:sweepdetail}
\end{table}

\begin{table}[h]
\centering\small
\begin{tabular}{rrrrrr}
\toprule
ToMe patches & MN patches & ToMe ms & MN ms & ToMe GiB & MN GiB \\
\midrule
1536 & 1538 & 147.30 & 178.17 & 2.793 & 3.619 \\
1151 & 1152 & 129.81 & 155.25 & 2.764 & 3.564 \\
922 & 923 & 122.44 & 143.02 & 2.746 & 3.721 \\
769 & 770 & 118.17 & 136.59 & 2.738 & 3.719 \\
575 & 576 & 106.45 & 129.38 & 2.723 & 3.557 \\
\bottomrule
\end{tabular}

\caption{Separate 384-pixel synthetic latency experiment: batch 64, H20-3e,
fp16 autocast, random initialization, \mn{} radius 5.143, and ToMe
proportional attention disabled. Patch counts exclude the class token here.
Memory is peak \emph{reserved} GiB, unlike the historical 224-pixel record,
which reports peak allocated memory.
Token-counting hooks remain active during timing.}
\label{tab:latency384}
\end{table}

At the approximately half-size operating point, ToMe processes 1,151 patches
and \mn{} 1,152, taking 129.810 and 155.255\,ms respectively. These times do not
prove a joint ordering for the accuracy sweep's ToMe setting. A common
harness should use exactly 1,152 patches with proportional attention enabled,
verify the checkpoint and preprocessing, and remove shape hooks before
timing. An earlier latency benchmark silently clamped several requested
ToMe budgets to the same achieved count; those obsolete cells are excluded.

Separate 384-pixel radius-3 benchmarks remain distinct from the radius-5.143
sweep. Speed percentages at different batch sizes, radii or harness settings
should not be averaged into an unexplained range. Peak reserved memory and
peak allocated memory also measure different quantities.

\section{Evidence and Reproduction Boundaries}

The artifact subset includes unmodified per-run summaries, resolved arguments
with host paths removed, launch-protocol extracts, sanitized successful sweep
records, audited benchmark tables, and source hashes. The CPU rebuild script
recomputes best/final metrics, resolves sweep records by timestamp, checks
numeric joins, and regenerates tables and figures. It never launches
training. Source and claim ledgers retain their literature retrieval dates;
the collaborator-packet audit was refreshed on September 21.
Company policy prevents ImageNet weight
transfer; independent ImageNet checkpoint re-evaluation is unavailable. The
September 23 packet additionally completes the two Stage-7 diagnostics.
The refreshed audit matches 3,354 of 3,560 summary rows to delivered EMA logs,
leaving 206 rows without corresponding log records
(Section~\ref{sec:setup}). Those rows are not discarded, but retain their
summary-only evidence status.

\section{DeiT-B/16 is not a usable scale anchor}
\label{app:deitb}

A 300-epoch DeiT-B/16 run at 224 pixels completes under the same global-batch
and seed protocol as the small-scale campaign, with drop-path 0.1 inherited
from the DeiT-S configurations. Table~\ref{tab:deitb} reports the peak and
final EMA scores. Train loss continues to fall while evaluation loss rises
after the peak at epoch 179. The final top-1, 78.566\%, is below our
DeiT-S/16 result of 79.692\% at the same 196-token budget. Published DeiT-B
300-epoch training uses drop-path 0.2. We therefore do not append this run
to the DeiT-S/8 table, and we do not launch a paired MergeNet-B comparison
from these weights.

\begin{table}[h]
\centering\small
\begin{tabular}{lrrr}
\toprule
Epoch & Train loss & Eval loss & EMA top-1 (\%) \\
\midrule
150 & 2.968 & 0.889 & 80.060 \\
179 (peak) & 2.808 & 0.895 & 80.442 \\
299 (final) & 2.348 & 1.095 & 78.566 \\
\bottomrule
\end{tabular}

\caption{DeiT-B/16, 300 epochs, drop-path 0.1. Peak and final scores must be
reported together; the peak is not a converged result.}
\label{tab:deitb}
\end{table}

\section{A Bounded Layout Optimization Experiment}
\label{app:layout}

We additionally test a single LocalAttention module on an A100-SXM4-80GB.
The original module constructs query/key/value tensors in
batch--head--token--channel order, then transposes and copies them for the
FlashAttention wrapper. An isolated variant constructs
batch--token--head--channel tensors directly and adjusts the global class-query
attention views and final reshape accordingly. It does not alter the
released model, load a checkpoint, or apply an optimizer update.

The test uses PyTorch 2.6.0+cu124 and FlashAttention 2.7.4.post1, width 384,
six heads, window radius 16, zero attention dropout, float32 parameters and
inputs, and fp16 autocast. Four batch-2 parity cases cover sequence lengths
785 and 2,305, global class attention on/off, and query/key normalization.
Outputs and input gradients are bitwise equal in these cases. Parameter
gradients are also equal except for the normalization case, whose maximum
relative $L_2$ error is $5.77\times10^{-7}$. This coverage does not establish
all-shape or dropout/RNG-preserving equivalence for a complete model.

\begin{table}[h]
\centering\small
\begin{tabular}{rrlrrr}
\toprule
Batch & Tokens & Mode & Original ms & Layout ms & Reduction range\\
\midrule
64 & 785 & Forward & 1.5600 & 1.0448 & 13.50--34.03\%\\
64 & 785 & Forward + backward & 5.1590 & 4.3999 & 14.68--25.64\%\\
16 & 2305 & Forward & 1.0318 & 0.9369 & 6.88--36.68\%\\
16 & 2305 & Forward + backward & 5.3740 & 4.9176 & 5.45--11.83\%\\
\bottomrule
\end{tabular}
\caption{Single-module A100 measurements. Times are medians of three round
medians, each with five warmups and twenty CUDA-event samples. The final
column is the range of paired-round time reductions, not a confidence
interval. Tokens include the class token. Backward includes a squared-output
loss and gradient clearing, with no optimizer.}
\label{tab:layout}
\end{table}

Measurement order alternates by round, and the harness requires an initially
empty device, establishes its CUDA process identity, then rejects a changed
process set at subsequent checks. All iteration samples and snapshots are
retained. The observed reductions support this local engineering direction,
but the round variation is appreciable. Both module instances are resident,
so their recorded memory peaks are not independent per-model memory
measurements. These results do not predict the full model's speedup, validate
an HBM-traffic estimate, transfer H20 timings to A100, or resolve the separate
ToMe attention-policy mismatch.
